\documentclass[12pt]{article}
\usepackage{amsmath, amssymb, amsthm}
\usepackage[margin=1in]{geometry}
\title{Higher Math HW3}
\author{Anatolii Larkin}

\begin{document}

\maketitle

\begin{enumerate}

\item Consider a simplified fragment of the genetic code. Let
\[
C = \{\text{AAA}, \text{AAG}, \text{UUU}, \text{UUC}, \text{GGG}\}, \quad P = \{\text{Lys}, \text{Phe}, \text{Gly}, \text{Trp}\}.
\]

Define \(f : C \to P\) using the table:
\[
\begin{array}{c|ccccc}
\text{Codon} & \text{AAA} & \text{AAG} & \text{UUU} & \text{UUC} & \text{GGG} \\
\hline
f(\text{Codon}) & \text{Lys} & \text{Lys} & \text{Phe} & \text{Phe} & \text{Gly}
\end{array}
\]

\begin{enumerate}
    \item Determine the image \(f(C)\).

    $im(f(C)) = \{f(x): x \in C\} = \{Lys, Phe, Gly\}$
    
    \item Is \(f\) injective? Explain.

    f is injective if $f(x_1) = f(x_2) \iff x_1 = x_2$.
    However $f(AAA) = f(AAG) = Lys$. Thus, f is not injective

    
    \item Is \(f\) surjective onto \(P\)? Explain.

    f is not surjective as $Trp \in P$, but there is not such $c \in C$ $that f(c) = Trp$
    
    \item Find \(f^{-1}(\{\text{Lys}\})\)

    $f^-^1(\{Lys\}) = \{AAA, AAG\}
    \item Would \(f : C \to f(C)\) be surjective? What changed?

    Since $C = \{AAA,AAG,UUU,UUC,GGG\}$ and $f(C) = \{Lys,Phe, Gly\}$, mapping $f: C \Rightarrow f(C)$ is surjective as for $\forall y \in f(C) \exists x \in C | f(x) = y$

\end{enumerate}


\item Let \(f, g : \mathbb{R} \to \mathbb{R}\) be defined by \(f(x) = 2x - 1\) and \(g(x) = x^2\).
    \begin{enumerate} 
        \item Find \((g \circ f)(x)\).

        $(g \circ f)(x) = g(f(x)) = (2x-1)^2
        
        \item Find \((f \circ g)(x)\).

        $(f \circ g)(x) = f(g(x)) = 2x^2-1
        
        \item Evaluate both compositions at \(x = 2\).

        $(g \circ f)(2) = 3^2 = 9

        $(f \circ g)(2) = 8 - 1 = 7$
        
        \item Does \(f \circ g = g \circ f\) hold? Explain.

        Assume \(f \circ g = g \circ f\).
        Then $2x^2 =1 = (2x-1)^2

        $x^2 - 2x + 1 = 0$

        $(x-1)^2 = 0$

        Thus, \(f \circ g = g \circ f\) holds iff $x = 1$. Thus it doesn't hold in general case
    \end{enumerate}
    
    \item A laboratory assigns samples to sequencing batches and batches to quality-control (QC) statuses. Let \(S = \{s_1, s_2, s_3, s_4\}\), \(B = \{b_1, b_2, b_3\}\) and \(Q = \{\text{pass}, \text{review}\}\). Define \(f : S \to B\) and \(g : B \to Q\) by
    \begin{align*}
        f(s_1) &= b_1, & f(s_2) &= b_1, & f(s_3) &= b_2, & f(s_4) &= b_3, \\
        g(b_1) &= \text{pass}, & g(b_2) &= \text{review}, & g(b_3) &= \text{pass}.
    \end{align*}
    \begin{enumerate}
        \item Construct the mapping \(g \circ f\) (give its values at all samples).

        $g\circf(s) = g(f(s)) : S \Rightarrow Q$
        $g(f(s_1)) = pass$
        $g(f(s_2)) = pass$
        $g(f(s_3)) = review$
        $g(f(s_4)) = pass$
        
        \item State its domain and codomain.

        $D(g(f(s)) =$ \{s_1,s_2,s_3,s_4\}
        
        $CD(g(f(s)) = \{pass, review\}$
        
        \item Is \(g \circ f\) injective? Is it surjective? Explain.

        $(g \circ f)(s)$ is not injective, as $g(f(s_1)) = g(f(s_2)) = pass, s_1 \not= s_2$
        $(g \circ f)(s)$ is surjective, as $\forall y \in CD(g(f(x)) \exists x \in D(g(f(x)) | g(f(x)) = y$
        
        \item Explain why \(f \circ g\) is not defined for these mappings.

        Assume mapping $f \circ g$ is defined. Then codomain Q of g must be a subset of domain S of f. However, $Q \not\subseteq S$. Contradiction found. Thus, mapping $f \circ g$ is not defined



    \end{enumerate}

    \item Work with the following mappings.
    \begin{enumerate}
        \item For \(f : \mathbb{R} \to \mathbb{R}\), \(f(x) = 3x-2\), find \(f^{-1}(x)\) and verify both identities \(f^{-1}(f(x)) = x\) and \(f(f^{-1}(x)) = x\).

        $f^-^1(x) = \frac{x+2}{3}$

        $f^-^1(f(x)) = \frac{3x-2+2}{3} = x$

        $f(f^-^1(x)) = 3(\frac{x+2}{3}) - 2 = x +2-2 = x$

        
        \item Explain why \(f : \mathbb{R} \to \mathbb{R}\), \(f(x) = x^2\), has no inverse mapping.

        Consider $x_1 = 2, x_2 = -2$. Thus $f(-2) = f(2) = 4, x_1 \not = x_2$. Thus, $f(x) = x^2$ is not injective. Thus, it is not bijective. Thus, its inverse mapping is undefined

        \item Consider instead \(f : [0, +\infty) \to [0, +\infty)\), \(f(x) = x^2\). Does an inverse exist? Find it if possible.

        Assume $f(x_1) = f(x_2)$. Then $x_1^2 = x_2^2$. Since $D(f(x)) = [0, +\infty), x_1 = x_2$. Thus, mapping is injective.

        Mapping $f(x)$ is also surjective as $\forall f(x) \in [0, +\infty) \exists x \in [0, +\infty) | x^2 = f(x)$

        Since mapping is surjective and injective, its bijective. Thus, its reverse function exists

        $f^-^1(x) = \sqrt{x}$

        $f^-^1(f(x)) = \sqrt{x^2} = x$

        $f(f^-^1(x)) = (\sqrt{x})^2 = x$

        
        \item Find the inverse of the mapping \(1 \mapsto 3, 2 \mapsto 1, 3 \mapsto 4, 4 \mapsto 2\) from \(\{1, 2, 3, 4\}\) to itself.

        The inverse mapping of such mapping is defined as follows:

        $1 \Rightarrow 2$
        
        $2 \Rightarrow 4$
        
        $3 \Rightarrow 1$
        
        $4 \Rightarrow 3$
        
    \end{enumerate}

    \item Let \(f : A \to B\) and \(g : B \to C\).
    \begin{enumerate}
        \item Prove that if \(f\) and \(g\) are injective, then \(g \circ f\) is injective.

        Assume $g \circ f$ is not injective

        Thus, $g(f(x_1)) = g(f(x_2)); x_1 \not = x2$

        Since g is injective, $g(f(x_1)) = g(f(x_2)), f(x_1) = f(x_2)$

        Therefore, $f(x_1) = f(x_2)$

        Since f is injective,
        $f(x_1) = f(x_2), x_1 = x_2$ \Rightarrow contradiction.

        Thus if f and g are injective, $g \circ f$ is injective
        
        \item Prove that if \(f\) and \(g\) are surjective, then \(g \circ f\) is surjective.

        Suppose $z \in C$. Since g is surjective, $\forall z \in C \exists y \in B | g(y) = z$. Since f is surjective, $\forall y \in B \exists x \in A | f(x) = y$. Combined, $\forall z \in C \exists x \in A | g(f(x)) = z$. Thus, if g and f are surjective, $g \circ f$ is surjective

        

        
    \end{enumerate}

    \end{enumerate}



    
\end{enumerate}




































\end{enumerate}
\end{document}